%% kapa: exam=1 round=34 year=2023 period=2 form=A subject=appraisal_law subject_ko=감정평가관계법규 seq=30 number=30
%% source: Q-Net 감정평가사 기출문제 (cst00309) artlSeq=5211789 file=raw/1008/5211789/2023년도 제34회 감정평가사 1차시험 2교시.pdf
%% conversion: kordoc 4.24.51 -> tools/gen1.py (auto); review=proofread-sonnet; flags=none
\begin{question}{30}
공간정보의 구축 및 관리 등에 관한 법령상 축척변경에 관한 설명으로 옳지 않은 것은?
\begin{choices}
  \item 하나의 지번부여지역에 서로 다른 축척의 지적도가 있는 경우 그 지역의 축척을 변경할 수 있다.
  \item 축척변경을 하려면 축척변경 시행지역의 토지소유자 2분의 1 이상의 동의를 받아야 한다.
  \item 합병하려는 토지가 축척이 다른 지적도에 각각 등록되어 있어 축척변경을 하는 경우 시·도지사 또는 대도시 시장의 승인을 요하지 않는다.
  \item 도시개발사업의 시행지역에 있는 토지로서 그 사업 시행에서 제외된 토지의 축척변경은 축척변경위원회의 의결을 요하지 않는다.
  \item 지적소관청은 축척변경 승인을 받았을 때에는 지체 없이 축척변경의 시행에 관한 세부계획을 20일 이상 공고하여야 한다.
\end{choices}
\end{question}
\answer{2}
